\chapter{Discussion}\label{ch:discussion}

The preceding chapter reported what we measured. This chapter asks what those measurements
\emph{mean}. We are deliberately careful here, because the central findings of this study are
largely null --- under a fair comparison, our from-scratch State-Optimization predictive coding
(PC) network neither beats nor is beaten by backpropagation (BP) at MNIST/FashionMNIST scale ---
and a null result is precisely the kind of result that invites overinterpretation in both
directions. We therefore separate four distinct interpretive questions and treat each on its own
terms: whether PC's reputed advantages are real (\Cref{sec:disc-real}), why a single methodological
checklist explains the literature's disagreement and transfers beyond PC
(\Cref{sec:disc-checklist}), what the systems measurements teach us about \emph{where the compute
actually goes} in PC inference (\Cref{sec:disc-systems}), what a generative PC genuinely buys that a
discriminative one does not (\Cref{sec:disc-generative}), and how an adversarial-verification habit
caught a framing error before it reached the page (\Cref{sec:disc-adversarial}). We close with an
explicit accounting of threats to the validity of the inferences drawn here
(\Cref{sec:disc-threats}). Throughout, the substantive enumeration of what this study \emph{cannot}
claim is deferred to \Cref{ch:limitations}; the present chapter is about interpretation, not
caveats.

\section{Are Predictive Coding's Advantages Real?}\label{sec:disc-real}

Predictive coding is advertised in the literature as a local-learning alternative to
backpropagation that, beyond its biological plausibility, is empirically \emph{better} along
several axes: it is said to approximate BP's gradients while needing fewer samples, to forget less
in continual-learning settings, to be more robust, and to do all of this with a single settling
mechanism that also doubles as a generative model. Our fair-comparison study
(\Cref{sec:exp-pcbp,sec:exp-song}) found none of the discriminative-side advantages to survive once
the comparison is made honestly. The headline verdict is parity: at matched loss, with
per-method-tuned learning rates and bootstrap confidence intervals over multiple seeds, PC and BP
are statistically indistinguishable on bulk accuracy, the confidence intervals overlap, and the
directional residue --- slightly worse noise robustness for PC, a hair's-breadth edge in the
faithful interference replication --- never exceeds one standard deviation.

\subsection{What ``parity'' does and does not mean}

It is worth stating the verdict with care, because ``PC $\approx$ BP'' is easy to misread. We are
\emph{not} claiming that PC and BP are the same algorithm. They are demonstrably not: PC settles the
free states to a (near-)equilibrium before touching the weights, whereas BP computes a single
backward pass and updates immediately, and the two procedures produce genuinely different
intermediate activity (this is the substance of Song et al.'s ``prospective configuration''
argument, \citealp{song2024prospective}). What we are claiming is narrower and more defensible:
along the \emph{outcome} metrics we measured --- test accuracy, noise robustness, sample efficiency,
and forgetting in continual learning --- the distinction between the two learning rules does not
manifest as a measurable performance gap once confounds are removed. The two rules arrive at
comparable solutions by different routes. That a different mechanism reaches the same place is itself
informative; it does not license the stronger claim that the mechanisms are equivalent, only the
claim that, at this scale and on these tasks, the mechanism does not buy a performance advantage.

Equally, parity is a \emph{scoped} claim, not a universal one. Our measurements live at
MNIST/FashionMNIST scale with shallow (two-to-five-layer) networks. The honest statement is that at
this scale, against a fairly tuned baseline, no PC advantage exceeds noise --- not that no advantage
can exist anywhere. The regimes in which PC is most strongly theorized to differ --- deep networks
where the prospective-configuration argument compounds across layers, genuine online single-example
streams, reinforcement learning --- are exactly the regimes our setup does not stress. This is the
correct frontier to defer to (\Cref{ch:limitations}), and we are careful not to overclaim past it.

\subsection{Why the literature disagrees}

If PC's advantages are not real at this scale, why does so much of the literature report them? Our
answer, and one of the transferable contributions of this study, is that the inconsistency is not
mysterious: naive comparisons trip over a small number of recurring confounds, and each confound
manufactures a spurious PC advantage in a predictable direction. We identified three (developed in
\Cref{sec:disc-checklist}), and in our own pre-confound runs all three fired. Our first
PC-versus-BP table showed an apparent BP \emph{win} on bulk accuracy and an apparent PC \emph{win}
on continual learning; the adversarial methodology review showed both to be artifacts. The bulk gap
was almost entirely the loss function --- cross-entropy versus mean-squared error to a one-hot
target --- not the learning rule, since the loss effect alone (BP-CE minus BP-MSE) accounts for the
bulk of a naive ``BP beats PC'' gap. The continual-learning gap was a double artifact: comparing
against a cross-entropy baseline (which forgets more because softmax produces large, aggressive
output updates) \emph{and} reading the forgetting metric without controlling for the fact that PC
learns each task less well in the first place, so it has mechanically less to forget.

This reframing reconciles what looks like a contradiction in the background literature. On the one
hand, a line of theoretical work shows that PC \emph{approximates} BP under restrictive conditions
--- the fixed-prediction regime of \citet{whittington2017approximation}, the arbitrary-graph result
of \citet{millidge2022backprop}, and the limiting analyses of \citet{innocenti2026limits}. On the
other hand, \citet{song2024prospective} argue that \emph{without} those constraints, energy-based
networks follow a distinct prospective-configuration dynamics that yields advantages. Both can be
true: PC collapses toward BP in the small-step, near-feedforward-initialization limit, and departs
from it when the settling is allowed to run to genuine equilibrium. Our finding sits on top of this:
even in the unconstrained, fully-settled regime where PC \emph{should} differ most, the difference
does not translate into a measurable advantage over a \emph{fairly tuned, loss-matched} BP at our
scale. The literature's disagreement is, on this reading, less about PC's true behaviour and more
about how carefully the BP baseline was constructed.

\section{The Confound Checklist as a Transferable Contribution}\label{sec:disc-checklist}

The most reusable product of this study is not the parity verdict itself but the diagnostic
machinery that produced it. We distilled the failure modes of naive PC-versus-BP comparisons into a
three-item checklist. We argue that this checklist is the genuine contribution of the methodological
strand of the thesis, because it generalizes well beyond PC: it applies to \emph{any} comparison
between a biologically-motivated learning rule and backpropagation, and indeed to many
``alternative-versus-baseline'' empirical claims in machine learning.

\subsection{The three confounds}

\paragraph{Confound 1: the loss function.}
The single largest hidden variable is the objective being minimized. A PC network that clamps a
one-hot target onto its linear output state and minimizes squared prediction error is minimizing
mean-squared error (MSE); a BP baseline that minimizes \texttt{CrossEntropyLoss} on logits is
optimizing a different loss surface entirely. These are not interchangeable: softmax cross-entropy
is well known to outperform MSE-to-one-hot on classification. Comparing PC-MSE against BP-CE
therefore conflates the learning rule with the loss, and because the loss effect is large
(\Cref{sec:exp-pcbp}), most of a naive gap is the objective, not the mechanism. The fix is to add a
matched BP-MSE arm so that the valid like-for-like comparison is PC-MSE versus BP-MSE, while the
difference BP-CE minus BP-MSE quantifies, and quarantines, the pure loss effect.

\paragraph{Confound 2: the stability--plasticity trade-off.}
In continual learning, ``forgets less'' is not the same as ``retains better.'' Backward transfer
measures the drop in accuracy relative to the freshly-learned diagonal; a model that writes weaker,
more diffuse task representations has a lower diagonal and therefore mechanically less to lose,
yielding a less-negative backward-transfer score that flatters it. A forgetting claim is only
defensible when read jointly against two controls: the per-task learn-accuracy (the diagonal, i.e.
how well each task was acquired) and the retained absolute accuracy on the first task after all
tasks (the sharpest single control, because it sidesteps the diagonal entirely). If PC ``forgets
less'' only because it learns less, the apparent advantage is a relabelling of underfitting.

\paragraph{Confound 3: the untuned baseline.}
A comparison is only as honest as its weakest baseline. We found that an untuned BP-MSE baseline
could sit stuck at chance on a class-incremental split, manufacturing a false PC ``win'' that a
learn-accuracy gate immediately exposed. The discipline this enforces is twofold: the baseline's
learning rate must be tuned on the same objective and validation protocol as the method under test,
and PC's own state and weight learning rates ($\eta_x$ and $\eta_w$) must be tuned \emph{jointly},
because they are coupled by a stability frontier --- sweeping $\eta_w$ alone, at a fixed and
too-large $\eta_x$, drives PC into divergence and systematically understates it.

\subsection{Why the checklist transfers}

None of these three confounds is specific to predictive coding. The loss confound bites whenever two
methods are compared under different objectives; the plasticity confound bites whenever a retention
metric is read without a learning control; the baseline confound bites whenever the comparison point
is left untuned. The checklist is thus a lightweight, reusable instrument for reviewers and authors
in an empirically confounded subfield: before believing a ``method X beats backpropagation'' claim,
ask whether the loss was matched, whether the retention metric was controlled for plasticity, and
whether the baseline was actually tuned to learn. We offer it in that spirit --- as a portable
artifact of the methodology, more durable than any single number in \Cref{ch:experiments}.

\section{The Systems Insight: Launch Overhead and the GEMM Trade-off}\label{sec:disc-systems}

The systems strand of the thesis produced the project's only unambiguous ``better'' result --- an
engineering win --- and it carries an interpretation worth drawing out independently of the speedup
figures. The measurement, not the speedup, is the contribution.

\subsection{PC inference is launch-overhead-bound, and that is a framework tax}

The headline systems observation (\Cref{sec:exp-kernel}) is that the PyTorch State-Optimization
backend spends roughly half a millisecond per settling step \emph{independent of batch size}, with
the GPU only partially utilized. The mechanistic evidence is direct: a launch-count profile shows
the PyTorch backend issuing thirty-one kernel launches per settling step (scaling linearly with the
number of settling steps $T$), whereas the fused kernel issues exactly one regardless of $T$. The
correct interpretation is that, at MNIST scale, PC inference is \emph{not} compute-bound; it is
bound by kernel-launch overhead --- the cost of dispatching many tiny operations from the host. This
is a recognized ``framework tax'' on workloads built from many small operations, and we are careful
to claim neither the observation nor the remedy as novel in general. The fused, persistent,
states-resident single-launch technique is established prior art, originating in cuDNN's persistent
RNN and Baidu's \texttt{persistent-rnn} \citep{diamos2016persistent} and reappearing in recent LLM
megakernels. What \emph{is} new, to our knowledge, is the transfer of that technique to the PC
settling loop and the measurement that pins the bottleneck: an audit of eleven open PC libraries and
the entire Equilibrium-Propagation family found none shipping a hand-written CUDA settling kernel
(\Cref{ch:related}; the audit is the substance of the novelty claim). The defensible statement is
therefore narrow --- ``first hand-written fused CUDA settling kernel \emph{for PC}, to our
knowledge'' --- and we hold it at that width deliberately.

\subsection{The architectural trade-off bounds the win}

The more interesting systems insight is \emph{why} the kernel wins only in a bounded regime. At
small and medium batch, fusing the $T$-step loop into one launch eliminates the launch overhead and
delivers a real speedup. But at large batch the bottleneck shifts from launch overhead to raw
compute, and here the design that wins at small batch actively works against you. A fused-resident
kernel holds the layer states on-chip across all $T$ steps; that residency consumes precisely the
register and shared-memory budget that a fast, register-blocked general matrix-multiply (GEMM) needs
to approach the hardware's floating-point peak. The fused-resident design and a cuBLAS-class GEMM are
thus \emph{architecturally opposed}: you cannot have both maximal residency and maximal GEMM
throughput in the same kernel, so at large batch the vendor library's register-blocked GEMM wins and
our naive in-kernel GEMM cannot keep up. This is not merely ``not yet optimized''; it is a structural
trade-off, and recognizing it as such is itself a reportable finding. Batch-tiling and caching the
clamped input activation push the crossover outward, but they do not abolish the trade-off.

\subsection{Why this matters for PC specifically}

The trade-off is not a defect for PC's purposes, because the regime where the kernel wins ---
small-to-medium batch --- is exactly the regime PC most naturally inhabits. PC is biologically
motivated as an online or near-online learner, and the prospective-configuration advantages are
theorized to appear in small-batch and single-example streams. The large-batch compute regime, where
cuBLAS is correctly the right tool, is the regime \emph{least} characteristic of PC. The practical
upshot is that the win lands where it is useful: when we ran the kernel as the engine of our own
faithful Song-exact study at batch~32 (\Cref{sec:exp-song}), it delivered a meaningful end-to-end
acceleration at per-seed-identical results, and its generalization to arbitrary depth (validated to
roughly $10^{-6}$ across depths two through five) means it is not a single-topology curiosity. This
is what ties the systems contribution to the science: the kernel is the measured, validated engine of
a real experiment, not a microbenchmark presented in isolation.

\section{What a Generative PC Buys}\label{sec:disc-generative}

One of PC's most cited selling points is the ``one mechanism, many tasks'' property: the same
settling rule that classifies should also generate, inpaint, and detect anomalies, with only the
clamping pattern changing. Our results (\Cref{sec:exp-gen}) make a sharp distinction here that we
believe is the correct interpretation, and it is easy to get wrong.

\subsection{The orientation, not the mechanism, is decisive}

A \emph{discriminatively} trained PC --- the standard supervised setup that clamps input and label
and settles the hidden states to learn an input-to-label map --- cannot generate. Clamping a label
and settling the input yields noise, not a recognizable digit, and residual energy fails to separate
in-distribution from out-of-distribution inputs. This is architectural, not a tuning failure: a
discriminatively trained network has learned a map from images to labels, not a generative model of
images, so settling ``backwards'' lands on no learned image manifold. The remedy is to flip the
orientation and train \emph{generatively} (label to image, clamping label and image while the hidden
states settle). With that single change, the identical settling rule that classified now produces
recognizable per-class prototypes, reconstructs occluded halves, and separates uniform-noise
out-of-distribution inputs with an anomaly AUC of $1.0$. The lesson is that the celebrated
``one mechanism, many tasks'' property is real, but it is a property of a \emph{generatively trained}
PC, not of the discriminative network most classification benchmarks use. The mechanism is the same;
the training objective is what unlocks the multi-task capability.

\subsection{Honest reading of the capability}

We are careful not to oversell what the generative PC delivers. With a one-hot label input there are
at most ten distinct outputs, so the ``generation'' is a set of class \emph{prototypes} --- the
per-class mean image --- not a diverse sample from a learned distribution; genuine sampling diversity
would require a continuous latent and a prior, as in a variational autoencoder. Inpainting is
prototype-driven and therefore blurry rather than pixel-sharp. And the anomaly result, while clean,
is measured against uniform noise, which is an easy out-of-distribution test; a near-OOD probe such
as FashionMNIST against MNIST, or rotated digits, would be a substantially harder and more
informative test. The defensible interpretation is therefore: a generative PC genuinely exhibits the
multi-task property from one settling rule, and the perfect anomaly separation is a real energy-based
signature (the model assigns high residual energy to inputs it cannot generate); but the qualitative
outputs are prototypes, and the OOD test is gentle. This is a demonstration of the property, scoped
honestly, not a competitive generative model.

\section{Methodology: Adversarial Verification}\label{sec:disc-adversarial}

A recurring theme of this study is that the most dangerous errors in an empirically confounded
subfield are not coding bugs --- those are caught by tests and gradient checks --- but framing
errors: subtle misreadings of what a prior result claims, or what one's own measurement isolates.
These are dangerous precisely because they survive a working pipeline. We adopted a lightweight
safeguard against them, and we argue it is reusable enough to be a methodological contribution in its
own right.

\subsection{The three-lens review and the error it caught}

Before committing a finding to the writeup, we subjected it to an adversarial review along three
independent lenses: a \emph{literature} lens (does this faithfully represent what the cited work
claims?), a \emph{statistics} lens (do the intervals, seed counts, and tests support the strength of
the claim?), and an \emph{alternative-explanations} lens (what confound could produce this result
without the mechanism we are attributing it to?). Applied to our Song-exact alternating replication
(\Cref{sec:exp-song}), this review caught a genuine framing error before it reached the page. Our
first draft attributed a learning-rate-robustness claim to Song's interference figure; the literature
lens established that the figure in question is a \emph{less-interference} claim measured with the
learning rate optimized \emph{per method}, while the learning-rate-robustness claim lives in a
\emph{different} figure. The correction was substantive: it changed what we tuned (we tune the
learning rate per method, as the interference claim requires) and what we report (we add the
worse-of-the-two-tasks accuracy as the interference tell, alongside the mean). Had the error
survived, we would have run and reported the wrong experiment while believing we were replicating
Song.

\subsection{Why this is a reusable safeguard, not a one-off}

The same three-lens discipline repeatedly corrected our own over-eager conclusions: the statistics
lens caught a confidence-interval mislabelling (68\% intervals reported as 95\%), and the
alternative-explanations lens is exactly what surfaced the three confounds of
\Cref{sec:disc-checklist}. We present adversarial verification not as a heroic one-time audit but as
a cheap, repeatable habit: for any empirical claim, ask whether it misreads the literature, whether
the statistics carry it, and whether a confound explains it just as well. In a subfield where
``method X beats backprop'' claims are easy to make and hard to reproduce, this is a low-cost,
high-yield form of self-skepticism. It is, in a sense, the methodological complement to the confound
checklist: the checklist enumerates the known traps, and the three-lens review is the procedure for
finding the ones not yet on the list.

\section{Threats to Validity}\label{sec:disc-threats}

Interpreting a null result responsibly requires being explicit about what could undermine the
inferences above. We organize the threats along the standard internal/external/construct axes and
state, for each, how it was mitigated and what residual risk remains. The substantive enumeration of
the study's scope limits --- scale, depth, prototype-only generation, the bounded kernel regime ---
belongs to \Cref{ch:limitations}; here we restrict ourselves to threats to the \emph{validity of the
conclusions} drawn in this chapter.

\subsection{Internal validity}

Internal validity concerns whether the measured differences (or non-differences) are caused by the
variable we intend to manipulate --- the learning rule --- rather than by an uncontrolled nuisance.
This is the axis the three confounds attack directly, and our principal mitigation is the
fair-comparison protocol itself: the BP baseline clones the PC network's initialization and uses the
identical forward function, so only the learning rule differs; the loss is matched via the BP-MSE
arm; learning rates are tuned per method on a held-out validation split and reported on test; and
PC's $\eta_x$ and $\eta_w$ are tuned jointly to respect the stability frontier. The residual internal
threats are honestly acknowledged: our headline used a modest seed budget with bootstrap intervals,
and although we raised the faithful replication to ten seeds --- which hardened rather than
overturned the verdict --- a non-significant directional effect at this seed count cannot be ruled to
be exactly zero, only ruled to be within noise. The learning-rate grids are finite, so we cannot
exclude that a finer grid would shift a point estimate, though the joint-tuning frontier we mapped
makes a large unmodelled improvement unlikely. We also note that the autonomous hyperparameter search
(\Cref{sec:exp-search}) independently rediscovered the $\eta_x \leftrightarrow \eta_w$ interaction,
which serves as an internal consistency check on the tuning.

\subsection{External validity}

External validity concerns how far the conclusions generalize. Here the threats are real and we do
not minimize them. All measurements are at MNIST/FashionMNIST scale with two-to-five-layer networks;
the parity verdict is scoped to that regime, and we explicitly decline to extend it to deep networks,
genuine online streams, or larger datasets, which is where PC is most theorized to differ. The
faithful Song replication, even at its most faithful, still differs from the original in scale and
seed budget, so our null does not globally refute Song --- the honest claim is the scoped one. The
single-GPU benchmarking (one RTX~3080~Ti) means the systems crossover points are hardware-specific;
the \emph{shape} of the launch-overhead-versus-compute trade-off is general, but the exact batch
sizes at which the kernel wins or loses would move on other hardware. These threats are mitigated by
scoping every claim to its measured regime rather than eliminated, and they motivate the future-work
directions of \Cref{ch:limitations}.

\subsection{Construct validity}

Construct validity concerns whether the metrics we chose actually measure the constructs we care
about. The risk is most acute in continual learning, where ``forgetting'' is a construct that a
single backward-transfer number measures poorly --- which is the entire point of the plasticity
confound. Our mitigation is to report forgetting as a triple (backward transfer, learn-accuracy, and
retained first-task accuracy) so that the construct is triangulated rather than reduced to one
potentially misleading scalar, and we added an EWC positive control that forgets less by
construction, confirming the harness measures forgetting correctly. For the anomaly-detection
capability, AUC against uniform noise is a valid but weak operationalization of ``out-of-distribution
detection,'' and we flag the easy-OOD caveat rather than letting the AUC of $1.0$ stand unqualified.
A subtler construct issue is that PC and BP differ in their \emph{inference} procedure --- PC settles
the output at test time, BP reads a single forward pass --- so noise-robustness and continual
comparisons partly measure inference dynamics, not only plasticity; we treat this as a documented
feature of PC rather than a flaw, but flag it so the reader knows the forward path is not perfectly
matched across arms. The recurring mitigation across all three axes is the same one that defines the
study: state every claim with its interval, its seed count, and its regime, and let the sober framing
do the work.
